\subsection{La direction sélectionnée par le registre dodécaphasique}
\label{exc:k-direccion}

Le registre \(K\) possède une fonction géométrique qui complète son évaluation rationnelle et sa lecture des supports. Une fois construite la représentation des douze positions, ses composantes déterminent une direction dans l'espace des variations de moyenne nulle. La réduction \(12\to11\) élimine le mode uniforme ; la réduction suivante \(11\to10\) élimine la direction que sélectionne le registre lui-même. Les deux opérations ont donc des antécédents différents.

Pour déclarer complètement la coordinatisation utilisée, considérons le groupe
\(A_5\) des permutations paires de cinq objets et le sous-groupe
\(C_5=\langle(1\,2\,3\,4\,5)\rangle\). Les douze classes à gauche
\(r_pC_5\) sont ordonnées au moyen de ces représentants, écrits sous forme de
listes d'images :
\[
\begin{array}{c|c@{\qquad}c|c}
p&r_p&p&r_p\\ \hline
1&(4,3,5,2,1)&7&(5,4,3,2,1)\\
2&(4,2,1,3,5)&8&(1,3,4,2,5)\\
3&(1,2,4,5,3)&9&(4,2,3,5,1)\\
4&(2,5,3,4,1)&10&(3,2,5,4,1)\\
5&(4,3,1,5,2)&11&(4,5,2,3,1)\\
6&(5,1,2,3,4)&12&(5,3,2,4,1).
\end{array}
\]
L'action par multiplication à gauche définit une représentation
orthogonale \(\rho:A_5\to O(\RR^{12})\) :
\(\rho(a)e_p=e_q\) lorsque \(ar_pC_5=r_qC_5\).
Cette coordinatisation de représentation agit sur le registre déjà obtenu ;
les caractères du groupe ne sont pas utilisés pour choisir ses douze valeurs.

Soient \(\chi_3\) les valeurs du caractère tridimensionnel
\[
\begin{array}{c|ccccc}
\text{classe}&1&2^2&3&5_a&5_b\\ \hline
\chi_3&3&-1&0&(1+\sqrt5)/2&(1-\sqrt5)/2.
\end{array}
\]
La classe \(5_a\) contient le générateur de \(C_5\) présenté et \(5_b\) contient son carré. Cette convention distingue les deux secteurs tridimensionnels conjugués de la représentation. Nous définissons
\begin{equation}
 P_3=\frac3{60}\sum_{a\in A_5}\chi_3(a^{-1})\rho(a),
 \qquad P_{11}=I_{12}-\frac1{12}J_{12}.
 \label{eq:k-proyectores-incidencia}
\end{equation}
Les deux matrices sont déterminées par des sommes finies dans
\(\QQ(\sqrt5)\). L'apparition de ce corps quadratique appartient
à la réalisation de la représentation, postérieure à la génération
régionale de \(\varphi\).

\begin{lemma}[Projection du registre sur le secteur tridimensionnel]
\label{lem:k-sector-no-nulo}
On a
\[
 P_3=P_3^{\mathsf T}=P_3^2,\quad
 P_3\mathbf1=0,\quad \operatorname{tr}P_3=3,
 \qquad P_3P_{11}=P_{11}P_3=P_3.
\]
Pour le registre \eqref{exc:k}, le vecteur
\begin{equation}
 u_K=P_3P_{11}K
 \label{eq:k-vector-direccional}
\end{equation}
est non nul. Sa norme exacte est
\begin{equation}
 \lVert u_K\rVert^2
 =\frac{6638585+2275584\sqrt5}{20}>0.
 \label{eq:k-norma-direccion}
\end{equation}
\end{lemma}
\begin{proof}
La somme de caractère dans \eqref{eq:k-proyectores-incidencia} est le
projecteur isotypique. Le caractère réel et la substitution \(a\mapsto a^{-1}\)
donnent la symétrie. La multiplicité du secteur dans l'action sur
\(A_5/C_5\) est la dimension de ses vecteurs fixes par \(C_5\),
\[
 \frac15\left(3+2\frac{1+\sqrt5}{2}
                 +2\frac{1-\sqrt5}{2}\right)=1.
\]
Le rang est donc trois. Son orthogonalité au caractère trivial donne \(P_3\mathbf1=0\) et les deux compositions avec \(P_{11}\). On peut également vérifier ces identités en multipliant les matrices de la somme finie.

Enfin, \(\sum K_p=6263\) et la substitution des douze
composantes dans cette même somme donnent
\[
 \lVert u_K\rVert^2=K^{\mathsf T}P_3K
 =\frac1{20}\sum_{a\in A_5}
                \chi_3(a^{-1})K^{\mathsf T}\rho(a)K
 =\frac{6638585+2275584\sqrt5}{20}.
\]
Il s'agit d'une évaluation exacte de soixante termes spécifiés,
reproduite par le certificat joint avec une arithmétique rationnelle
quadratique. Les deux coefficients positifs prouvent la non-nullité.
\end{proof}

\begin{definition}[Direction dodécaphasique et complément transversal]
La direction sélectionnée par \(K\) dans la coordinatisation déclarée est
\(\ell_K=\RR u_K\). Son complément transversal dans
\(V_{11}=\mathbf1^\perp\) est
\[
 V_{10}(K)=V_{11}\cap u_K^\perp.
\]
\end{definition}

\begin{proposition}[Réduction orthogonale déterminée par \(K\)]
\label{prop:k-reduccion-diez}
L'opérateur
\begin{equation}
 P_{10}(K)=P_{11}
       -\frac{u_Ku_K^{\mathsf T}}{\lVert u_K\rVert^2}
 \label{eq:k-proyector-diez}
\end{equation}
est le projecteur orthogonal de rang dix sur \(V_{10}(K)\). On a
\[
 V_{11}=\ell_K\mathbin{\oplus^{\perp}}V_{10}(K),
 \qquad P_{10}P_{11}=P_{10},\quad P_{10}u_K=0.
\]
\end{proposition}
\begin{proof}
Soit \(E_K=u_Ku_K^{\mathsf T}/\lVert u_K\rVert^2\). La non-nullité démontrée garantit qu'il est défini. Il est symétrique, \(E_K^2=E_K\), de rang un et, puisque \(u_K\in V_{11}\), \(P_{11}E_K=E_KP_{11}=E_K\). D'où
\[
 (P_{11}-E_K)^2=P_{11}-2E_K+E_K=P_{11}-E_K.
\]
Son image est \(V_{11}\cap u_K^\perp\) et sa trace est
\(11-1=10\). La décomposition et les identités restantes
s'obtiennent par projection sur les deux termes orthogonaux.
\end{proof}

Ce résultat précise une fonction de \(K\) que sa dénomination de registre n'épuisait pas : outre la représentation réversible de données orientées et la sélection de la tétrade \(B_K\), il détermine une polarisation linéaire. La direction dépend du registre complet par l'intermédiaire de \(P_3K\), et non uniquement des quatre points de son support. L'ordre registre, représentation, direction et réduction dimensionnelle a ainsi été maintenu (proposition~\ref{prop:k-reduccion-diez}).

La transformation possède une loi de covariance précise. Si l'on change
simultanément la coordinatisation par une matrice orthogonale de permutation \(S\),
\(K'=SK\) et \(P'_3=SP_3S^{-1}\), alors
\[
 u_{K'}'=Su_K,\qquad P'_{10}(K')=SP_{10}(K)S^{-1}.
\]
D'autre part, remplacer \(u_K\) par \(a u_K\), avec \(a\ne0\),
ne change pas \(E_K\). La direction retient une polarisation et ne
permet pas de retrouver à elle seule l'amplitude de \(u_K\) ni le
registre intégral. La réversibilité de \(\kappa\leftrightarrow K\)
demeure dans son domaine propre.

La réalisation géométrique établie ici est une réduction d'espaces à produit scalaire. Son emploi ultérieur comme direction compacte exige de transporter cette droite vers la réalisation physique correspondante. La section \ref{exc:pantallas-dualidad} expose l'application algébrique qui maintient coordonnées cette réduction et l'écran réticulaire.
